% Shared preamble for both drivers: the monograph and the extractable Part I.
% Conventions follow paper/workhouse_publication_edition_rev5_2026-08-30.tex
% so that material can move between editions without renaming symbols.

\usepackage[T1]{fontenc}
\usepackage[utf8]{inputenc}
\usepackage{lmodern}
\usepackage[margin=1in,headheight=14pt]{geometry}
\usepackage{amsmath,amssymb,amsthm,mathtools}
\usepackage{booktabs,array,longtable,enumitem}
\usepackage{xcolor,microtype,xurl}
\usepackage{fancyhdr}
\usepackage[hidelinks]{hyperref}

\hypersetup{colorlinks=true,
  linkcolor=blue!45!black, citecolor=blue!45!black, urlcolor=blue!45!black}

\setlength{\emergencystretch}{3em}
\setlength{\parskip}{2pt}
\setlist{itemsep=2pt,topsep=4pt}
\numberwithin{equation}{section}

\theoremstyle{plain}
\newtheorem{theorem}{Theorem}[section]
\newtheorem{proposition}[theorem]{Proposition}
\newtheorem{lemma}[theorem]{Lemma}
\newtheorem{corollary}[theorem]{Corollary}
\newtheorem{conjecture}[theorem]{Conjecture}
\theoremstyle{definition}
\newtheorem{definition}[theorem]{Definition}
\newtheorem{problem}{Problem}
\theoremstyle{remark}
\newtheorem{remark}[theorem]{Remark}
\newtheorem{scopenote}[theorem]{Scope}

% ---------------------------------------------------------------- notation
\newcommand{\R}{\mathbb{R}}
\newcommand{\C}{\mathbb{C}}
\newcommand{\Q}{\mathbb{Q}}
\newcommand{\Z}{\mathbb{Z}}
\newcommand{\N}{\mathbb{N}}
\newcommand{\SU}{\mathrm{SU}}
\newcommand{\Tr}{\operatorname{Tr}}
\newcommand{\spec}{\operatorname{spec}}
\newcommand{\Ric}{\operatorname{Ric}}
\newcommand{\im}{\operatorname{im}}
\newcommand{\rank}{\operatorname{rank}}
\newcommand{\norm}[1]{\left\lVert #1 \right\rVert}
\newcommand{\abs}[1]{\left\lvert #1 \right\rvert}
\newcommand{\ip}[2]{\left\langle #1, #2 \right\rangle}
\newcommand{\Hil}{\mathcal{H}}
\newcommand{\Alg}{\mathcal{A}}
\newcommand{\gp}{\mathfrak{g}}
\newcommand{\Hph}{\Hil_{\mathrm{phys}}}

% ------------------------------------------------- verification-tier badges
% The badge records how a statement is machine-certified in the repository.
% It is INDEPENDENT of mathematical status: a T3 badge on a proven theorem
% means the proof is analytic and no machine check covers the whole
% statement.  Conflating the two is the error this notation exists to stop.
\definecolor{tierzero}{RGB}{20,90,40}
\definecolor{tierone}{RGB}{25,70,140}
\definecolor{tiertwo}{RGB}{130,90,10}
\definecolor{tierthree}{RGB}{100,100,105}

\newcommand{\tierbadge}[2]{%
  \textcolor{#1}{\footnotesize\textsf{[#2]}}}
\newcommand{\Tzero}{\tierbadge{tierzero}{T0 Lean}}
\newcommand{\Tone}{\tierbadge{tierone}{T1 exact}}
\newcommand{\Ttwo}{\tierbadge{tiertwo}{T2 numeric}}
\newcommand{\Tthree}{\tierbadge{tierthree}{T3 analytic}}

% Graph identity of a statement, so a reader or an agent can run
% `workhouse why <id>` on anything asserted here.
\newcommand{\gid}[1]{\par\nobreak\smallskip\noindent
  {\footnotesize\textsf{Graph id:} \nolinkurl{#1}}\par\smallskip}
\newcommand{\gidc}[2]{\par\nobreak\smallskip\noindent
  {\footnotesize\textsf{Graph id:} \nolinkurl{#1}\quad
   \textsf{Reproduce:} \texttt{#2}}\par\smallskip}
\newcommand{\srcpath}[1]{\nolinkurl{#1}}

% An explicit, visible statement of what a theorem does NOT say.  Used
% wherever the corpus records a scope limit, which is most places.
\newcommand{\notasserted}[1]{\par\smallskip\noindent
  {\small\textbf{Does not assert.} #1}\par\smallskip}

% The obligation register's problem headers.
\newcommand{\obligation}[2]{\par\medskip\noindent
  \textbf{#1}\ \ \textit{#2}\par\smallskip}
